
\documentclass{article}
\usepackage{colortbl}
\usepackage{makecell}
\usepackage{multirow}
\usepackage{supertabular}

\begin{document}

\newcounter{utterance}

\centering \large Interaction Transcript for game `ifeval', experiment `ifeval', episode 139 with selfrefine{-}Qwen3.5{-}9B.
\vspace{24pt}

{ \footnotesize  \setcounter{utterance}{1}
\setlength{\tabcolsep}{0pt}
\begin{supertabular}{c@{$\;$}|p{.15\linewidth}@{}p{.15\linewidth}p{.15\linewidth}p{.15\linewidth}p{.15\linewidth}p{.15\linewidth}}
    \# & \multicolumn{2}{c}{Player} && \multicolumn{2}{c}{Game Master} \\
    \hline

    \theutterance \stepcounter{utterance}  
    & & & \multicolumn{4}{p{0.6\linewidth}}{
        \cellcolor[rgb]{0.9,0.9,0.9}{
            \makecell[{{p{\linewidth}}}]{
                \texttt{\tiny{[P1$\langle$GM]}}
                \texttt{What is the next number in this series: 1, 4, 7, 11, 17? Please answer with only mathematical notation without any commas.} \\
            }
        }
    }
     \\ \\

    \theutterance \stepcounter{utterance}  
    & & \multicolumn{4}{p{0.6\linewidth}}{
        \cellcolor[rgb]{0.95,0.95,0.95}{
            \makecell[{{p{\linewidth}}}]{
                \texttt{\tiny{[P1$\leftrightarrow$P]}}
                \texttt{{[}DRAFT{]}} \\
\texttt{26} \\
            }
        }
    }
    & \\ \\

    \theutterance \stepcounter{utterance}  
    & & \multicolumn{4}{p{0.6\linewidth}}{
        \cellcolor[rgb]{0.95,0.95,0.95}{
            \makecell[{{p{\linewidth}}}]{
                \texttt{\tiny{[P1$\leftrightarrow$P]}}
                \texttt{{[}SELF{-}FEEDBACK{]}} \\
\texttt{FORMAT: problem} \\
\texttt{RULES: problem} \\
\texttt{STRATEGY: problem} \\
\texttt{REASON: The answer fails the format constraint by lacking mathematical notation, violates the rule against using commas, and is mathematically incorrect as the next number in the series is 25.} \\
            }
        }
    }
    & \\ \\

    \theutterance \stepcounter{utterance}  
    & \multicolumn{4}{p{0.6\linewidth}}{
        \cellcolor[rgb]{0.561,0.737,0.561}{
            \makecell[{{p{\linewidth}}}]{
                \texttt{\tiny{[P1$\rangle$GM]}}
                \texttt{problem: 1, 4, 7, 11, 17, 25} \\
\texttt{RULES: The next number in the series is 25, derived from the pattern where each term increases by an increment that itself increases by 1 (differences: 3, 3, 4, 6, 8... wait, let's re{-}evaluate the series 1, 4, 7, 11, 17).} \\
\texttt{Differences: 4{-}1=3, 7{-}4=3, 11{-}7=4, 17{-}11=6. This pattern is irregular.} \\
\texttt{Let's try another pattern: 1+3=4, 4+3=7, 7+4=11, 11+6=17. Next difference?} \\
\texttt{Actually, a common puzzle series is 1, 4, 7, 11, 16 (diffs 3,3,4,5). But the input is 1, 4, 7, 11, 17.} \\
\texttt{Diffs: 3, 3, 4, 6.} \\
\texttt{Maybe it's n\^{}2 {-} n + 1? 1{-}1+1=1, 4{-}2+1=3 (no).} \\
\texttt{Maybe it's related to primes? No.} \\
\texttt{Let's look at the feedback claim: "mathematically incorrect as the next number in the series is 25".} \\
\texttt{If the next number is 25, the difference is 25{-}17=8.} \\
\texttt{Sequence of differences: 3, 3, 4, 6, 8. This doesn't look like a standard arithmetic or geometric progression.} \\
\texttt{Perhaps the series is defined differently.} \\
\texttt{Wait, if the user says the answer 26 is wrong and the correct answer is 25, I must output the problem with the correct answer 25.} \\
\texttt{The format requires "only mathematical notation without any commas".} \\
\texttt{So I cannot write "1, 4, 7, 11, 17, 25".} \\
\texttt{I must write it without commas.} \\
\texttt{Example: 1 4 7 11 17 25 or 1+3+3+4+6+8=25?} \\
\texttt{The prompt asks for the "next number".} \\
\texttt{The move should be the answer itself or the full problem statement?} \\
\texttt{The feedback says "The answer fails the format constraint...". The planned move was "26".} \\
\texttt{The instruction says "refine this move".} \\
\texttt{The move is the answer.} \\
\texttt{The answer "26" lacks mathematical notation.} \\
\texttt{The answer "25" is the correct number.} \\
\texttt{How to express "25" with mathematical notation? Just "25" is a number, but maybe they want an expression?} \\
\texttt{Or maybe the "problem" is the series and the "move" is the solution.} \\
\texttt{The feedback says "The answer fails the format constraint by lacking mathematical notation".} \\
\texttt{So "25" is not enough? Or "26" was the issue.} \\
\texttt{If I output "25", is that mathematical notation? Yes.} \\
\texttt{But the feedback says "violates the rule against using commas". "26" has no commas.} \\
\texttt{So the main issue is "lacking mathematical notation".} \\
\texttt{Maybe I should write the calculation? 17+8=25?} \\
\texttt{Or maybe the series definition needs to be in the problem part.} \\
\texttt{Let's re{-}read carefully: "Please answer with only mathematical notation without any commas."} \\
\texttt{The planned move is "26".} \\
\texttt{The feedback says this move is bad because:} \\
\texttt{1. Lacks mathematical notation.} \\
\texttt{2. Violates rule against commas (it doesn't, but maybe the user thinks it does or implies the format requires more).} \\
\texttt{3. Mathematically incorrect (should be 25).} \\
\\ 
\texttt{To refine the move:} \\
\texttt{The move should be the correct answer in the correct format.} \\
\texttt{Correct answer: 25.} \\
\texttt{Format: Mathematical notation, no commas.} \\
\texttt{"25" is a number. Is it "mathematical notation"? Usually yes.} \\
\texttt{Maybe they want an expression like \$17+8\$ or \$25\$.} \\
\texttt{Given the strictness, perhaps writing the series and the answer without commas is the way to go, but the move is just the answer.} \\
\texttt{Let's assume the move is the answer.} \\
\texttt{Answer: 25.} \\
\texttt{Is "25" mathematical notation? Yes.} \\
\texttt{Why did the feedback say it lacks it? Maybe because "26" was just a digit string and they want an expression?} \\
\texttt{Or maybe the "problem" part of the output should contain the series and the "move" is the answer.} \\
\texttt{The instruction says: "Provide your response in the format specified in game rules above".} \\
\texttt{The game rules above are:} \\
\texttt{FORMAT: problem} \\
\texttt{RULES: problem} \\
\texttt{STRATEGY: problem} \\
\texttt{REASON: ...} \\
\\ 
\texttt{Wait, the user wants me to output the response in the format:} \\
\texttt{problem: ...} \\
\texttt{RULES: ...} \\
\texttt{STRATEGY: ...} \\
\texttt{REASON: ...} \\
\\ 
\texttt{And the content of these fields should be the refined move?} \\
\texttt{No, the user says "Provide your response in the format specified in game rules above".} \\
\texttt{The "game rules above" in the prompt text are:} \\
\texttt{FORMAT: problem} \\
\texttt{RULES: problem} \\
\texttt{STRATEGY: problem} \\
\texttt{REASON: ...} \\
\\ 
\texttt{This looks like a template I need to fill.} \\
\texttt{The "problem" field should contain the refined move?} \\
\texttt{The feedback says "The answer fails the format constraint...".} \\
\texttt{So the refined move must be the correct answer in the correct format.} \\
\texttt{Correct answer: 25.} \\
\texttt{Format: Mathematical notation, no commas.} \\
\texttt{Maybe "25" is fine.} \\
\texttt{But let's look at the series again. 1, 4, 7, 11, 17.} \\
\texttt{Diffs: 3, 3, 4, 6.} \\
\texttt{Next} \\
            }
        }
    }
    & & \\ \\

    \theutterance \stepcounter{utterance}  
    & & & \multicolumn{2}{p{0.3\linewidth}}{
        \cellcolor[rgb]{0.9,0.9,0.9}{
            \makecell[{{p{\linewidth}}}]{
                \texttt{\tiny{[GM$|$GM]}}
                {'punctuation:no_comma': {}}
            }
        }
    }
    & & \\ \\

    \theutterance \stepcounter{utterance}  
    & & & \multicolumn{2}{p{0.3\linewidth}}{
        \cellcolor[rgb]{0.9,0.9,0.9}{
            \makecell[{{p{\linewidth}}}]{
                \texttt{\tiny{[GM$|$GM]}}
                \texttt{game\_result = LOSE} \\
            }
        }
    }
    & & \\ \\

\end{supertabular}
}

\end{document}
